\documentclass[12pt]{article}
\usepackage{ae}
\usepackage[T1]{fontenc}
\usepackage{amsmath}
\usepackage{amsfonts}
\usepackage{lmodern}
\usepackage[lmargin=1.5cm, rmargin=2cm,tmargin=2cm,bmargin=2cm]{geometry}
\usepackage{graphicx}
\usepackage{xcolor}
\usepackage{booktabs}
\usepackage{array}
\usepackage{enumitem}
\usepackage{xurl}
\usepackage{hyperref}

\title{Group Project: Causal Inference Paper Exposition \\ Quantitative Methods for Humanities}
\author{}
\date{2026--27}

\setlist[itemize]{noitemsep, topsep=2pt}
\setlist[enumerate]{noitemsep, topsep=2pt}

\begin{document}
\maketitle

\section*{Purpose}

The goal of this project is to show that your group can read, understand, and explain a published paper that uses causal inference. 
You are not being asked to carry out a new causal analysis in R or to work with the assigned course datasets. 
Instead, your task is to study one paper carefully and give a clear, 
brief exposition of its research question, evidence, causal strategy, results, and limitations.

The project is designed to reward understanding rather than technical decoration. 
A good presentation should make the paper understandable to classmates who have taken this course but have not read the paper.

\section*{Basic Rules}

\begin{itemize}
    \item Groups may have at most \textbf{three students}.
    \item The project is worth \textbf{10\%} of the final grade.
    \item Each group will choose and analyze one \textbf{causal-inference paper} from the list of papers provided by the instructor.
    \item One group member should submit the slides and the short report as \textbf{PDF files} on Canvas.
    \item The deliverables are:
    \begin{itemize}
        \item a short oral presentation, supported by clear and informative slides,
        \item a short report introducing and explaining the paper.
    \end{itemize} 
    \item Every member of the group must understand the whole paper. The instructor may ask any member to answer any question.
\end{itemize}

\section*{Paper Selection}

Each group will either choose or be assigned a paper from the list provided in the \hyperref[sec:papers]{Papers section} at the end of this document. 

\subsection*{Proposing your own paper}

If your group wants to propose another paper, ask for approval before working on it. A suitable paper should:

\begin{itemize}
    \item study a social, political, cultural, historical, educational, economic, or policy question;
    \item make an explicit causal claim or evaluate a causal effect;
    \item identify a treatment, intervention, exposure, reform, event, or assignment mechanism;
    \item define at least one outcome variable;
    \item use evidence that can be explained with the causal inference concepts covered in class.
\end{itemize}

\section*{Presentation Format}

The presentation should be at most \textbf{10 minutes} (longer presentations will be penalized), followed by a short round of questions. 
As a guideline, use no more than \textbf{6 content slides}, plus one final slide for references if needed.

Slides should be clear and selective. Language should be precise and straightforward, not vague or ambiguous. 
Do not try to summarize every detail of the paper. Focus on the core causal argument and the evidence needed to understand it.

\section*{Short Report Format}

The short report should follow the same structure as the presentation. It must be at most \textbf{5 pages}, excluding references, and use a readable format: \textbf{12 pt font}, at least \textbf{2.5 cm margins}, and \textbf{1.5 line spacing}. Do not use dense formatting to fit more text.

\section*{Required Content}

Your presentation and report should cover the following points. Some papers may not contain all of them in the same way. 
If something is not applicable, say so briefly and explain why.

\subsection*{1. Paper, Problem, and Motivation}

Begin with the full citation of the paper. Explain:

\begin{itemize}
    \item what problem the paper studies;
    \item why the problem matters;
    \item what gap, debate, or puzzle the paper tries to address.
\end{itemize}

\subsection*{2. Research Question and Main Claim}

State the main research question in your own words. If possible, write it as a causal question. %: ``What is the effect of $T$ on $Y$?'' 
Then state the paper's main answer or claim.

Be precise. Avoid vague descriptions such as ``the paper studies inequality'' or ``the paper talks about education.'' Instead, identify what is being compared and what outcome is being measured.

\subsection*{3. Causal Framework}

Identify the main causal elements of the paper:

\begin{itemize}
    \item the unit of analysis;
    \item the treatment, intervention, or exposure;
    \item the outcome variable or variables;
    \item the comparison or control group;
    \item the causal effect or estimand, when the paper defines one;
    \item the counterfactual idea behind the paper's argument.
\end{itemize}

If the paper has several treatments or outcomes, focus on the most important one and mention the others only if needed and time allows.

\subsection*{4. Data and Descriptive Evidence}

Explain the data used in the paper:

\begin{itemize}
    \item the data source;
    \item the sample and time period;
    \item the key variables;
    \item any descriptive table, figure, or exploratory evidence that helps motivate the analysis.
\end{itemize}

You may reproduce or adapt one figure or table from the paper if it is central to the argument. If you do, cite the original table, figure, or page clearly.

\subsection*{5. Causal Inference Technique}

Explain the paper's identification strategy in words. 
Possible techniques include randomized experiments, natural experiments, instrumental variables, regression discontinuity, difference-in-differences, matching, panel fixed effects, or another design.

Your explanation should answer:

\begin{itemize}
    \item Why do the authors think their comparison can be interpreted causally?
    \item What assumption is most important for the causal interpretation?
    \item What could threaten that assumption?
    \item What checks, placebo tests, balance tests, robustness checks, or sensitivity analyses do the authors use, if any?
\end{itemize}

Use and relate to course concepts when relevant: treatment, outcome, control group, potential outcomes, selection bias, confounding, randomization, parallel trends, average treatment effect, internal and external validity.
If needed to understand the paper, you must research other concepts not explored in the course.

\subsection*{6. Results, Interpretation, and Uncertainty}

Present the main result. Explain the direction, size, and meaning of the estimated effect in plain language.

If the paper reports confidence intervals, standard errors, $p$-values, or hypothesis tests, explain what they add to the interpretation. Do not simply say that a result is ``significant.'' Explain what hypothesis is being tested and what the evidence suggests.

\subsection*{7. Conclusions and Limitations}

End by explaining:

\begin{itemize}
    \item what the authors conclude;
    \item whether you think the causal interpretation is convincing;
    \item the strongest limitation of the paper;
    \item what you learned about causal inference and its applications from reading the paper.
\end{itemize}

This section should include your group's judgment, not only a summary of the authors' conclusion.

% \section*{Use of AI}

% You may use AI tools to help understand difficult passages, define technical terms, improve slide wording, or generate practice questions. AI may also be useful for asking, ``What should I check in this identification strategy?'' or ``What does this robustness check mean?''

% However, AI summaries are not a substitute for reading the paper. Your group is responsible for every claim made in the slides and presentation. You should be able to point to the relevant page, table, figure, or equation in the paper for your main claims.

% \subsection*{Allowed Uses}

% \begin{itemize}
%     \item Asking for explanations of unfamiliar terminology.
%     \item Asking AI to quiz you on the paper after you have read it.
%     \item Improving the clarity or grammar of your slides.
%     \item Asking for feedback on a draft explanation written by your group.
% \end{itemize}

% \subsection*{Not Allowed}

% \begin{itemize}
%     \item Presenting an AI-generated summary without checking it against the paper.
%     \item Inventing details about data, methods, results, or limitations.
%     \item Hiding meaningful AI use.
%     \item Using AI to replace the group's own reading and causal reasoning.
% \end{itemize}

% \subsection*{AI Use Note}

% On the final slide, include a short AI use note with three items:

% \begin{enumerate}
%     \item Which AI tool(s), if any, you used.
%     \item What you used AI for.
%     \item One example of something AI suggested that your group checked, changed, or rejected.
% \end{enumerate}

% If you did not use AI, write: ``No AI tools were used for this project.''

% \section*{Oral Questions}

% After the presentation, the instructor may ask questions to any member of the group. Possible questions include:

% \begin{itemize}
%     \item What is the paper's main research question?
%     \item What is the treatment? What is the outcome?
%     \item What is the comparison group or counterfactual?
%     \item What causal inference technique does the paper use?
%     \item What is the key identification assumption?
%     \item What is the strongest threat to the causal interpretation?
%     \item What is the main result, in plain language?
%     \item What table, figure, or result in the paper is most important?
%     \item What did your group find difficult or unconvincing?
%     \item How did your group use AI, if at all?
% \end{itemize}

% \section*{Rubric}

% \begin{center}
% \begin{tabular}{p{5.1cm}p{1.4cm}p{7.0cm}}
% \toprule
% \textbf{Criterion} & \textbf{Points} & \textbf{What matters} \\
% \midrule
% Problem and research question & 1.5 & Clear explanation of the paper's problem, motivation, research question, and main claim. \\
% \addlinespace
% Causal framework & 2.0 & Correct identification of treatment, outcome, unit of analysis, comparison group, estimand, and counterfactual logic. \\
% \addlinespace
% Data and descriptive evidence & 1.5 & Accurate explanation of the data source, sample, key variables, and most relevant descriptive evidence. \\
% \addlinespace
% Causal method and assumptions & 2.0 & Clear explanation of the identification strategy, key assumptions, threats to validity, and checks used by the authors. \\
% \addlinespace
% Results and interpretation & 1.5 & Main results are explained in plain language, with careful interpretation of effect sizes, uncertainty, hypothesis tests, and limitations. \\
% \addlinespace
% Presentation, questions, and AI transparency & 1.5 & Slides are clear and selective, all members can answer basic questions, and AI use is declared honestly. \\
% \bottomrule
% \end{tabular}
% \end{center}

% \section*{Grade Caps}

% The following issues may cap the project grade, regardless of how polished the slides look:

% \begin{itemize}
%     \item If the group cannot identify the research question, treatment, outcome, or causal method, the project cannot receive more than 5/10.
%     \item If the presentation mainly repeats generic AI-style summaries and cannot be connected to specific pages, tables, or figures in the paper, the project cannot receive more than 6/10.
%     \item If the group misrepresents the paper's data, method, or results in a serious way, the project cannot receive more than 6/10.
%     \item If meaningful AI use is hidden, the project cannot receive more than 6/10.
%     \item If one or more members cannot answer basic questions about the paper, their individual project grade may be capped.
%     \item Fabricated sources, fabricated results, or serious academic dishonesty may receive 0/10.
% \end{itemize}

\section*{Submission Checklist}

Before submitting, check that your slides and report include:

\begin{itemize}
    \item the names of all group members;
    \item the full citation of the paper;
    \item the paper's research question and main claim;
    \item treatment, outcome, comparison group, and causal method;
    \item the main data source and one key table, figure, or result;
    \item a careful interpretation and one limitation;
    \item references to the relevant paper pages, tables, or figures;
\end{itemize}

\phantomsection
\section*{Papers}\label{sec:papers}

{\raggedright
\begin{enumerate}[leftmargin=*]

    \item Salinas, Paula, and Albert Sol\'e-Oll\'e. 2018. ``Partial Fiscal Decentralization Reforms and Educational Outcomes: A Difference-in-Differences Analysis for Spain.'' \emph{Journal of Urban Economics} 107: 31--46. \href{https://doi.org/10.1016/j.jue.2018.08.003}{doi: 10.1016/j.jue.2018.08.003}.

    \item Felfe, Christina, Natalia Nollenberger, and N\'uria Rodr\'iguez-Planas. 2015. ``Can{'}t Buy Mommy{'}s Love? Universal Childcare and Children{'}s Long-Term Cognitive Development.'' \emph{Journal of Population Economics} 28(2): 393--422. \href{https://doi.org/10.1007/s00148-014-0532-x}{doi: 10.1007/s00148-014-0532-x}.

    \item Brassiolo, Pablo. 2016. ``Domestic Violence and Divorce Law: When Divorce Threats Become Credible.'' \emph{Journal of Labor Economics} 34(2): 443--477. \href{https://doi.org/10.1086/683666}{doi: 10.1086/683666}.

    \item Fern\'andez-Kranz, Daniel, and N\'uria Rodr\'iguez-Planas. 2021. ``Too Family Friendly? The Consequences of Parent Part-Time Working Rights.'' \emph{Journal of Public Economics} 197: 104407. \href{https://doi.org/10.1016/j.jpubeco.2021.104407}{doi: 10.1016/j.jpubeco.2021.104407}.

    \item Garc\'ia-Hombrados, Jorge, Marta Mart\'inez-Matute, and Carmen Villa Llera. 2024. ``Specialised Courts and the Reporting of Intimate Partner Violence: Evidence from Spain.'' \emph{Journal of Public Economics} 239: 105243. \href{https://doi.org/10.1016/j.jpubeco.2024.105243}{doi: 10.1016/j.jpubeco.2024.105243}.

    \item Torun, Huzeyfe. 2019. ``Ex-Ante Labor Market Effects of Compulsory Military Service.'' \emph{Journal of Comparative Economics} 47(1): 90--110. \href{https://doi.org/10.1016/j.jce.2018.12.001}{doi: 10.1016/j.jce.2018.12.001}.

    \item Gorj\'on, Luc\'ia, David Mart\'inez de Lafuente, and Gonzalo Romero. 2024. ``Employment Effects of the Minimum Wage: Evidence from the Spanish 2019 Reform.'' \emph{SERIEs: Journal of the Spanish Economic Association} 15: 1--55. \href{https://doi.org/10.1007/s13209-023-00291-1}{doi: 10.1007/s13209-023-00291-1}.

    \item Arranz, Jos\'e M., and Carlos Garc\'ia-Serrano. 2025. ``Assessing the Impact of an Increase in the Minimum Wage on Household Income and Poverty.'' \emph{Social Science Research} 127: 103143. \href{https://doi.org/10.1016/j.ssresearch.2025.103143}{doi: 10.1016/j.ssresearch.2025.103143}.

    \item Garc\'ia-Vega, Mar\'ia, Richard Kneller, and Joel Stiebale. 2021. ``Labor Market Reform and Innovation: Evidence from Spain.'' \emph{Research Policy} 50(5): 104213. \href{https://doi.org/10.1016/j.respol.2021.104213}{doi: 10.1016/j.respol.2021.104213}.

    \item Gonz\'alez-Val, Rafael, and Miriam Marc\'en. 2024. ``Population Growth and Taxes: The Effect of Regional Differences in the Spanish Inheritance Tax.'' \emph{The Annals of Regional Science} 73: 135--163. \href{https://doi.org/10.1007/s00168-023-01253-y}{doi: 10.1007/s00168-023-01253-y}.

    \item Yoshimura, Yuji, Yusuke Kumakoshi, Yichun Fan, Sebastiano Milardo, Hideki Koizumi, Paolo Santi, Juan Murillo Arias, Siqi Zheng, and Carlo Ratti. 2022. ``Street Pedestrianization in Urban Districts: Economic Impacts in Spanish Cities.'' \emph{Cities} 120: 103468. \href{https://doi.org/10.1016/j.cities.2021.103468}{doi: 10.1016/j.cities.2021.103468}.

    \item Vald\'es, Manuel T., Mar C. Espadafor, and Risto Conte Keivabu. 2025. ``Can a Low Emission Zone Improve Academic Performance? Evidence from a Natural Experiment in the City of Madrid.'' \emph{Population and Environment} 47: 8. \href{https://doi.org/10.1007/s11111-025-00477-8}{doi: 10.1007/s11111-025-00477-8}.

    \item Pinilla, Jaime, Beatriz G. L\'opez-Valc\'arcel, and Miguel A. Negr\'in. 2019. ``Impact of the Spanish Smoke-Free Laws on Cigarette Sales, 2000--2015: Partial Bans on Smoking in Public Places Failed and Only a Total Tobacco Ban Worked.'' \emph{Health Economics, Policy and Law} 14(4): 536--552. \href{https://doi.org/10.1017/S1744133118000270}{doi: 10.1017/S1744133118000270}.

    \item Lid\'on-Moyano, Cristina, Marcela Fu, Montse Ballb\`e, Juan Carlos Mart\'in-S\'anchez, Nuria Matilla-Santander, Cristina Mart\'inez, Esteve Fern\'andez, and Jose M. Mart\'inez-S\'anchez. 2017. ``Impact of the Spanish Smoking Laws on Tobacco Consumption and Secondhand Smoke Exposure: A Longitudinal Population Study.'' \emph{Addictive Behaviors} 75: 30--35. \href{https://doi.org/10.1016/j.addbeh.2017.06.016}{doi: 10.1016/j.addbeh.2017.06.016}.

    \item Tassinari, Filippo. 2024. ``Low Emission Zones and Traffic Congestion: Evidence from Madrid Central.'' \emph{Transportation Research Part A: Policy and Practice} 185: 104099. \href{https://doi.org/10.1016/j.tra.2024.104099}{doi: 10.1016/j.tra.2024.104099}.

    \item Abadie, Alberto, and Javier Gardeazabal. 2003. ``The Economic Costs of Conflict: A Case Study of the Basque Country.'' \emph{American Economic Review} 93(1): 113--132. \href{https://doi.org/10.1257/000282803321455188}{doi: 10.1257/000282803321455188}.

    \item Gonz\'alez, Libertad. 2013. ``The Effect of a Universal Child Benefit on Conceptions, Abortions, and Early Maternal Labor Supply.'' \emph{American Economic Journal: Economic Policy} 5(3): 160--188. \href{https://doi.org/10.1257/pol.5.3.160}{doi: 10.1257/pol.5.3.160}.

\end{enumerate}
\par}

\end{document}
